\documentclass[conference]{IEEEtran}
\IEEEoverridecommandlockouts
\usepackage{amsmath,amssymb,amsthm}
\usepackage{cite}
\usepackage{algorithm}
\usepackage{algpseudocode}
\usepackage{microtype}
\usepackage{url}
\usepackage{tikz}
\usepackage{etoolbox}
\usepackage[hidelinks]{hyperref}

\newtheorem{theorem}{Theorem}
\newtheorem{lemma}{Lemma}
\newtheorem{corollary}{Corollary}
\newcommand{\dist}{\operatorname{dist}}
\newcommand{\E}{\mathbb{E}}
\newcommand{\Prb}{\mathbb{P}}
\newcommand{\cL}{\mathcal{L}}

\begin{document}

\title{Exact Near-Girth Cycle Counting in Tanner Graphs:\\
Algorithms and Hybrid Estimation}

\author{\IEEEauthorblockN{Arash Ahadi}
\IEEEauthorblockA{Tehran Institute for Advanced Studies (TeIAS), Khatam University, Iran\\
\texttt{aarash.ahadi.academic@gmail.com}}}

\maketitle

\begin{abstract}
Short cycles in Tanner graphs of LDPC codes create dependencies among messages in iterative
decoding and can degrade decoding performance. Consequently, determining
their number is important in the analysis and design of LDPC codes. We consider exact and randomized counting of short cycles in a realized Tanner
graph.

Let a Tanner graph have girth $g\ge6$, order $M$, and maximum degree $\Delta$.
For every even $g\le T\le M$, we give a deterministic algorithm that
simultaneously computes $N_g,N_{g+2},\ldots,N_T$ in
\[
O\!\left(gM\Delta^{\,T-g/2+1}\right)
\]
time.

As a complementary practical component for $(d_v,d_c)$-regular graphs, we
analyze an unbiased uncapped inverse-sampling estimator for $N_g$ based on
non-backtracking walks and combine it with the exact routine through an
independent pilot test and a linear cap. Increasing the stopping target gives
explicit finite-sample control of both relative error and failure probability.
Experiments validate exact counts beyond the near-girth range and show
target- and graph-dependent runtime gains and reversals against our
Karimi--Banihashemi implementation. Randomized estimation is effective
when girth cycles are sufficiently abundant.
\end{abstract}

\begin{IEEEkeywords}
LDPC codes, Tanner graph, cycle counting, girth, non-backtracking walk, Chernoff's bound, BFS.
\end{IEEEkeywords}

\section{Introduction}
Low-density parity-check (LDPC) codes~\cite{gallager1962} are among the most
widely used families of modern error-correcting codes. Their sparse
parity-check structure is naturally represented by a Tanner
graph~\cite{tanner1981},
a bipartite graph whose variable nodes represent code symbols and whose check
nodes represent parity constraints. Iterative decoding operates locally on
this graph, and short cycles cause messages to become correlated after only a
few iterations, which can degrade decoding performance. Consequently, the
number and distribution of short cycles are important
structural parameters in LDPC-code analysis and design
\cite{halford2006,karimi2013,blake2018,dehghan2019}.

Write $M=|V(G)|$ and let $\Delta$ be the maximum degree;
$T$ denotes the largest requested cycle length.
For general bipartite graphs, Halford and Chugg~\cite{halford2006} gave an
early short-cycle counting algorithm. Karimi and
Banihashemi~\cite{karimi2013} developed a distributed message-passing method
that counts all lengths $g,g+2,\ldots,2g-2$ in $O(g|E|^2)$ time; Li, Lin,
and Abdel-Ghaffar~\cite{li2015} subsequently reduced its computational cost
by a constant factor. Dehghan and Banihashemi~\cite{dehghan2020} later gave
a BFS-based method running in $O(M^2\Delta)$ for $g$- and $(g+2)$-cycles
and $O(M^2\Delta^2)$ for $(g+4)$-cycles, together with hardness results for
longer lengths. More recently, Chen and He~\cite{chen2025fmpa}
developed a fast message-passing algorithm (FMPA) to reduce message-processing work; its published bound
remains quadratic in $|E|$ for fixed $g$ and $T$.
Our local search complements these approaches.
For fixed $g$, $\Delta$, and $T$, our bound is linear in $M$, versus the
published quadratic KB bound; this compares upper bounds, not guaranteed
runtime behavior. Our bound is most attractive for small degrees and
short targets, and grows exponentially with $T$.
Section~\ref{sec:experiments} reports both gains and reversals.

Related work also includes general cycle-counting
algorithms~\cite{alon1997}, message-passing methods for counting
and enumerating cycles in bipartite graphs~\cite{yang2017},
and parallel algorithms specialized to induced
$6$-cycles~\cite{niu2025}.

Spectral methods use adjacency or directed-edge spectra and degree data
\cite{blake2018,dehghan2019,dehghan2020directed}.
QC methods exploit circulant shifts, protographs, or lifting structure
\cite{fossorier2004,karimi2012qc,gomez2023}.
Our exact method instead uses local BFS and residual-path enumeration
on the realized graph, allowing irregular graphs without these structures;
our randomized analysis assumes $(d_v,d_c)$-regularity.

Since $4$-cycles are undesirable in iterative LDPC
decoding, we assume girth $g\ge6$ throughout.

In Section~\ref{sec:exact}, we count all cycles
of every even length $t$ satisfying
\begin{equation}\label{eq:range}
 g\le t\le M.
\end{equation}
More strongly, for any even $T$ in this range, all counts
$N_g,N_{g+2},\ldots,N_T$ are obtained simultaneously in
\begin{equation}\label{eq:maincomplexity}
 O\!\left(gM\Delta^{\,T-g/2+1}\right),
\end{equation}
The key structural observation is that, for each chosen root, the girth
condition forces the two initial segments of every cycle through that
root to follow unique paths in its local BFS tree. These tree paths
provide a fixed part of the cycle, reducing the search to a residual
path connecting their endpoints. We exploit this decomposition to count
all requested lengths simultaneously, including lengths beyond the
near-girth range. The analysis accounts for both residual-path
enumeration and the disjointness checks needed to ensure that each
completed cycle is simple, yielding the stated complexity bound.

When girth cycles are sufficiently abundant, sampling can provide
an efficient alternative to exact counting. In Section~\ref{sec:random}, the closure probability of sampled length-$g$
non-backtracking walks determines $N_g$: each closure traverses a girth cycle.
Non-backtracking walks are established tools in spectral
theory~\cite{bass1992};
we use this girth-specific counting identity in a hybrid estimator.

A short pilot then decides whether closures are too rare for efficient
relative sampling. Closure-sparse instances are routed to the exact routine,
whereas cycle-rich instances use fresh independent inverse samples. We prove explicit finite-sample
routing bounds and exponentially decaying relative-error bounds in the
stopping target.

\section{Exact Cycle Counting}\label{sec:exact}
Let $G=(U\cup W,E)$ be a simple bipartite graph of known girth $g\ge6$
and maximum degree $\Delta$. Girth computation, if needed, is preprocessing.
Write $\cL_i(v)=\{x:\dist(v,x)=i\}$. The following elementary lemmas expose
the tree structure used by the exact algorithm.

\begin{lemma}[Tree ball]\label{lem:tree}
For every $v$, the subgraph induced by vertices at distance at most $g/2-1$
from $v$ is a tree. Thus each $x\in\cL_{g/2-1}(v)$ has a unique
path from $v$ to $x$, denoted by $P_v(x)$.
\end{lemma}
\begin{proof}
The BFS tree spans the ball, so any cycle in the induced ball contains a
non-tree edge $xy$. This edge together with the tree paths from $v$ to $x$
and $y$ contains a cycle of length at most
\[
\dist(v,x)+\dist(v,y)+1\le
\left(\frac{g}{2}-1\right)+\left(\frac{g}{2}-1\right)+1
=g-1,
\]
contradicting the assumption that the graph has girth $g$.
\end{proof}

Lemma~\ref{lem:tree} yields the saving in
\eqref{eq:maincomplexity}: two tree arms are fixed,
leaving only the residual path to search.

For $x\in\cL_{g/2-1}(v)$, let $b_v(x)$ be the first vertex after $v$ on
$P_v(x)$.  Fix any even $t$ with $g\le t\le M$ and set
\begin{equation}\label{eq:rdef}
 r=t-g+2.
\end{equation}
Then $r\ge2$; no upper bound on $r$ in terms of $g$ is assumed.

\begin{lemma}[Rooted decomposition]\label{lem:decomp}
Let $C$ be a $t$-cycle and $v\in C$.  Moving $g/2-1$ edges from $v$ in the
two directions on $C$ reaches two vertices in $\cL_{g/2-1}(v)$, called $x$ and $y$, in different
branches.  The selected arcs are $P_v(x)$ and $P_v(y)$, and the remaining
$x$--$y$ arc has length $r$.
\end{lemma}
\begin{proof}
Each of the two selected arcs from $v$ to $x$ and $y$ has length $g/2-1$.
If, say, $x$ had a shorter $v$--$x$ path, its union with the selected arc
would contain a cycle of length at most $g-3$, contradicting the girth
assumption. Thus $x,y\in\cL_{g/2-1}(v)$. Moreover, a second shortest path
to either endpoint, together with the corresponding selected arc, would
contain a cycle of length at most $2\left(g/2-1\right)=g-2$.
Hence the selected arcs are precisely the unique BFS-tree paths
$P_v(x)$ and $P_v(y)$.
They leave $v$ along the two different cycle edges, so they lie in different branches. Here, each branch of the BFS tree rooted at $v$ is the subtree rooted at one
of the children of $v$.
The total length of $P_v(x)$ and $P_v(y)$ is $g-2$, leaving exactly
$r=t-(g-2)$ edges.
\end{proof}



\begin{algorithm}[t]
\caption{Exact counting of $t$-cycles}\label{alg:exact}
\footnotesize
\begin{algorithmic}[1]
\Require Tanner graph $G=(U\cup W,E)$ with girth $g\ge6$; even
$t$ satisfying $g\le t\le M$
\State $r\gets t-g+2$, $C_t\gets0$; fix any total order $\prec$ on $V(G)$ to break symmetry
\ForAll{$v\in U$}
  \State build the BFS tree to depth $g/2-1$
  \ForAll{$x\in\cL_{g/2-1}(v)$}
    \State enumerate length-$r$ non-backtracking walks $Q$ from $x$ whose internal vertices avoid $P_v(x)$
    \Statex \hspace{\algorithmicindent} using DFS that marks the current path and rejects repeated vertices
    \ForAll{such generated paths $Q$}
      \State let $y$ be the endpoint of $Q$
      \If{$y\in\cL_{g/2-1}(v)$, $b_v(x)\ne b_v(y)$,
          $x\prec y$, and $Q$ internally avoids $P_v(y)$}
        \State $C_t\gets C_t+1$
      \EndIf
    \EndFor
  \EndFor
\EndFor
\State \Return $N_t=2C_t/t$
\end{algorithmic}
\end{algorithm}

\begin{theorem}\label{thm:exact}
Let $G$ be a Tanner graph with girth $g\ge6$ and maximum degree $\Delta$, and let
$g\le T\le M$ be any even integer. Then all cycle counts $N_g,N_{g+2},\ldots,N_T$
can be computed exactly in $O\!\left(gM\Delta^{\,T-g/2+1}\right)$
time. 
\end{theorem}
\begin{proof}
Fix an even $t\in\{g,g+2,\ldots,T\}$ and put $r=t-g+2$. By
Lemmas~\ref{lem:tree} and~\ref{lem:decomp}, every $t$-cycle through a root
$v\in U$ is split uniquely into two forced BFS-tree arms
$P_v(x),P_v(y)$, each of length $g/2-1$, together with the remaining
$x$--$y$ path of length $r$.

Conversely, every residual non-backtracking walk accepted by Algorithm~\ref{alg:exact}
is simple because in line~5 of Algorithm~\ref{alg:exact}, the DFS never extends its current path to a vertex
already marked on that path. Since its internal vertices avoid both tree
arms, and the two arms intersect only at $v$, their union is a simple
$t$-cycle. This argument applies to every $r$, including $r\ge g$.

For each $x$, mark $V(P_v(x))\setminus\{x\}$ in a separate forbidden
array in $O(g)$ time. These marks remain fixed during the search from
$x$. A second array marks the vertices on the current DFS path,
initially only $x$. DFS explores each neighbor that is neither
forbidden nor already on the current path, marking it before recursion
and unmarking it on return. Each neighbor check and mark update takes
$O(1)$ time. At depth $r$, DFS tests the endpoint and acceptance
conditions. Thus every admissible simple residual path is generated
exactly once from $x$. To check that $Q$ internally avoids $P_v(y)$
(line~8 of Algorithm~\ref{alg:exact}), scan
$V(P_v(y))\setminus\{y\}$ against the current-path marks, rejecting
the candidate if any scanned vertex is marked. This ensures that $Q$
meets the second tree arm only at $y$ and costs $O(g)$, as accounted
for below. These acceptance tests are performed before backtracking,
while all vertices of $Q$ are still marked.

Distinct accepted residual paths for the same $\{x,y\}$ yield distinct
cycles and must all be counted. 
Cycles formed solely by the union of residual $x$--$y$ paths are not
counted at this root, because every residual path avoids $v$, whereas each
counted cycle includes $v$ through the two tree arms.

The unique rooted decomposition and $x\prec y$ give a bijection between
accepted residual paths and $t$-cycles through $v$; the residual connection
need not be unique.

Every $t$-cycle contains exactly $t/2$ vertices of $U$, each of which
is used as a root. Thus the raw counter satisfies
\[
C_t=\frac{t}{2}N_t,
\]
and returning $2C_t/t$ yields the exact cycle count $N_t$.

To count all lengths through $T$ simultaneously, for each root $v$ and each
$x\in\cL_{g/2-1}(v)$ perform one non-backtracking search to maximum depth
$R=T-g+2$, forbidding $V(P_v(x))\setminus\{x\}$ and rejecting repeated
vertices by the same current-path marking. Whenever the current path
reaches depth $r=t-g+2$ and ends at a boundary vertex $y$ in a different
branch, with $x\prec y$, check that its internal vertices avoid $P_v(y)$
and update the counter for length $t$ if the check succeeds. Reaching a boundary
vertex does not stop the search: longer simple extensions may contribute
to larger target lengths. The cycle--path correspondence above applies
at every admissible depth.

\underline{Complexity:}
For each root $v\in U$, there are $O(\Delta^{g/2-1})$ choices of $x$.
From each $x$, DFS generates $O(\Delta^{T-g+2})$ partial paths.
Neighbor checks occur only below the maximum search depth;
their total number is also $O(\Delta^{T-g+2})$ per $x$.
The endpoint $y$ is determined by each generated path, so all possible
endpoints are already included in this bound. Checking whether a
candidate path avoids the second tree arm $P_v(y)$ internally costs
$O(g)$. Since there are $n=|U|$ roots, the total cost is
\[
O\!\left(ng\Delta^{g/2-1}\Delta^{T-g+2}\right)
=
O\!\left(ng\Delta^{T-g/2+1}\right).
\]
Since $n\le M$, this proves the stated bound.

Allocate the BFS and marking arrays once in $O(M)$ time.
Do not clear all $M$ entries for each new root. Instead, record which
entries are used and clear only those entries before reusing the arrays.
The number of entries cleared is at most the number used during the
search, so cleanup is included in the stated complexity bound.
We count each arithmetic operation on a cycle counter as one operation.
\end{proof}

\section{Accuracy-Controlled Hybrid $g$-Cycle Counting}
\label{sec:random}

When girth cycles are sufficiently abundant, sampling
can provide an efficient alternative to exact counting; see inequality (\ref{eq:rich-regime}).
When such cycles are sparse, exact counting may be preferable.

We combine the exact routine of Section~\ref{sec:exact} with inverse
sampling. The two procedures
are interleaved on one processor, receiving equal budgets of
elementary operations. The first to finish supplies the answer.

Assume that $G$ is $(d_v,d_c)$-regular, with $n=|U|$ variable
nodes and girth $g\ge6$. 
Selecting a non-backtracking walk means that, at each step, we choose a neighbor of the vertex selected in the previous step, excluding the vertex selected two steps earlier. In other words, no edge may be traversed in the reverse direction immediately after it has been traversed.

The number of length-$g$ non-backtracking
(non-backtracking) walks from each variable node is
\begin{equation}\label{eq:B}
 B=d_v(d_v-1)^{g/2-1}(d_c-1)^{g/2}.
\end{equation}
A sampling trial chooses a variable root uniformly and then a
uniform non-backtracking walk of length $g$ from that root. By the girth
assumption, a closed such walk traverses a simple $g$-cycle.
Each $g$-cycle contributes two orientations from each of its
$g/2$ variable vertices. Thus the closure probability is
\begin{equation}\label{eq:closure-prob}
 p=\frac{gN_g}{nB}.
\end{equation}

\subsection{Inverse sampling and accuracy selection}
For $N_g>0$, let $S_r$ be the sample on which the $r$th
closure occurs, $r\ge2$. Then
$S_r\sim\mathrm{NegBin}(r,p)$ and $\mathbb E[S_r]=r/p$.
The standard unbiased inverse-sampling estimator
of $p$ is $(r-1)/(S_r-1)$ for $r\ge2$
~\cite{mendo2009}. Using \eqref{eq:closure-prob},
our estimator of $N_g$ is
\begin{equation}\label{eq:inverse-est}
 \widehat N_g=\frac{nB}{g}\frac{r-1}{S_r-1}.
\end{equation}
It is unbiased without stopping by the exact procedure.
More importantly, the normalized mean absolute-error bound
of~\cite{mendo2009} implies
\begin{equation}\label{eq:mendo-bound}
 \mathbb E\!\left[\frac{|\widehat N_g-N_g|}{N_g}\right]
 \le a_r:=
 \frac{2(r-1)^{r-1}e^{-(r-1)}}{(r-1)!}.
\end{equation}
Given a requested MARE tolerance $\varepsilon>0$, choose
\begin{equation}\label{eq:r-eps}
 r=r(\varepsilon)=\min\{j\ge2:a_j\le\varepsilon\}.
\end{equation}
The target is independent of $N_g$; indeed
$a_r\sim\sqrt{2/(\pi r)}$, so $r=O(\varepsilon^{-2})$
as $\varepsilon\rightarrow0$.

For standard accuracy tolerances, the corresponding
values of $r(\varepsilon)$ can be precomputed and
retrieved in constant time. For arbitrary tolerances,
$r(\varepsilon)$ can be obtained in $O(r)$ numerical
operations using the recurrence
$a_{j+1}=a_j e^{-1}(j/(j-1))^{j-1}$, starting
from $a_2=2/e$.
\subsection{Single-processor interleaving}
Both procedures are made resumable. The exact routine stores
its BFS and DFS states, root index, path marks, and counters;
the sampler stores its current (possibly incomplete) walk,
trial count $S$, closure count $A$, and random-generator state.
Fix any constant quantum $Q\ge1$. The scheduler gives at most
$Q$ \emph{elementary operations} to each procedure in turn,
not one full iteration of each: a single exact-search
iteration and a sampling trial need not have equal cost.
A procedure may yield in the middle of a search or walk and
resume from that point. No physical clock or second core is
required.

\begin{algorithm}[t]
\caption{Interleaved hybrid $g$-cycle counting}
\label{alg:hybrid}
\footnotesize
\begin{algorithmic}[1]
\Require $(d_v,d_c)$-regular Tanner graph $G$ with $n$ variable vertices and girth
$g\ge6$, tolerance $\varepsilon>0$
\State Fix any constant $Q\ge1$
\State Compute $B$ from \eqref{eq:B}
\State $r\gets r(\varepsilon)$ from \eqref{eq:r-eps}
\State Initialize resumable $\mathsf{ExactGirth}(G)$
\State Initialize sampler with $S\gets0$, $A\gets0$
\While{true}
  \State Advance $\mathsf{ExactGirth}$ by at most $Q$ operations
  \If{exact counting finishes}
    \State \Return $N_g$
  \EndIf
  \State Advance sampling by at most $Q$ operations
  \State Update $S$ for completed walks and $C$ for closures
  \If{$C\ge r$}
 \State \Return $\frac{nB}{g}\frac{C-1}{S-1}$
  \EndIf
\EndWhile
\end{algorithmic}
\end{algorithm}
\begin{theorem}\label{thm:hybrid}
For $N_g>0$, Algorithm~\ref{alg:hybrid} returns
$\widetilde N_g$ satisfying
\[
 \mathbb E\!\left[
  \frac{|\widetilde N_g-N_g|}{N_g}\right]
 \le\varepsilon.
\]
For $N_g=0$ it returns zero exactly. Under unit-cost
arithmetic and $O(1)$ amortized suspension and resumption,
its work after initialization is
\[
 O\!\left(\min\{W_E,W_R\}+Q\right)
 =O(gn\omega+Q),
\]
where $W_E$ is the work of exact girth counting,
$W_R$ is the work of uncapped inverse sampling to its
$r$th closure (infinite when $N_g<r$), and
$\omega$ is the regular-degree factor in
Section~\ref{sec:exact} for $T=g$.
\end{theorem}
\begin{proof}
Couple the hybrid sampler to an infinite sequence of
independent trials. If sampling wins, the output equals
\eqref{eq:inverse-est}; otherwise it is the exact count.
Therefore its absolute error is no larger, sample by
sample, than that of the uncapped estimator.
Taking expectations and applying \eqref{eq:mendo-bound}
proves the MARE claim. If $N_g<r$, the exact routine finishes first.

With indexed adjacency lists, a completed non-backtracking trial takes
$O(g)$ elementary operations. Because each active routine
receives the same constant quantum $Q$ per round, the
first completion occurs within
$2\min\{W_E,W_R\}+O(Q)$ elementary operations, up to
constant-amortized scheduling overhead. By the regular-degree
bound of Section~\ref{sec:exact},
$W_E=O(gn\omega)$; the work bound follows.
\end{proof}

\noindent\textbf{When sampling is advantageous.}
Uncapped sampling needs an expected $r/p$ trials,
each costing $\Theta(g)$ elementary operations. Its expected
work, apart from initialization, is therefore
\begin{equation}\label{eq:random-work}
 \mathbb E[W_R]
 =\Theta\!\left(\frac{gr}{p}\right)
 =\Theta\!\left(\frac{rnB}{N_g}\right).
\end{equation}
If sampling wins, the exact routine imposes an upper cutoff
on $S_r$; hence its conditional expected number of trials
is at most $r/p$. The expected total work of the race obeys
\begin{equation}\label{eq:hybrid-expected-work}
 \mathbb E[W_{\rm hybrid}]
 =O\!\left(
 \min\left\{gn\omega,\frac{rnB}{N_g}\right\}+Q
 \right),\qquad N_g>0.
\end{equation}
Thus the \emph{displayed bounds} favor sampling when
\begin{equation}\label{eq:rich-regime}
 N_g\gg \frac{rB}{g\omega}
 =\Theta\!\left(\frac{B}{\varepsilon^2 g\omega}\right).
\end{equation}
and favor exact counting when $N_g$ is small compared with
this scale. Decreasing $\varepsilon$ increases
$r=\Theta(\varepsilon^{-2})$, shifting the comparison
toward exact counting. For fixed degrees, girth, and $\varepsilon$,
$\omega$, $B$, and $r$ are constants; if
$N_g=\Theta(n)$, expected sampling work is $O(1)$
apart from input setup, versus the $O(n)$ exact bound.
If $N_g=O(1)$, the sampling-work expression is
$\Theta(n)$, and exact counting can be preferable.
These comparisons concern asymptotic work bounds, not
finite-instance wall-clock speedups or a lower bound on
the exact procedure. The MARE guarantee is an expectation,
not a per-run error guarantee.

\section{Experimental Evaluation}\label{sec:experiments}
Single-threaded C++17 measurements use \texttt{g++ 14.2 -O3} on an
AMD EPYC 9V74 shared container. Exact/KB times are medians of seven runs;
hybrid times and outputs average 100 runs, including exact returns.
Graph construction and verification are untimed. Source code, graph generators, seeds, edge lists, and full results
for the comparison are available at
\url{https://github.com/aarashahadiacademic-1/tanner-cycle-counting}.

\subsection{Benchmark Construction for Comparison with KB
\cite{karimi2013}}
The graphs in Table~\ref{tab:comparison-exp} are synthetic sparse parity-check
matrices, not standardized codes or decoding benchmarks.
We favor small degrees to limit the branching of our local search.

The labels in the first column of Table~\ref{tab:comparison-exp}
identify the graph construction and, where applicable, its lift.
Bases A--C use point--line incidence $y=sx+b$ over $\mathbb F_Q$,
the finite field with $Q$ elements, with
$0\le x \le d_c-1$ and $0\le s\le d_v-1$.
Their parameters $(Q,d_v,d_c)$ are $(7,3,6)$, $(5,3,4)$, and
$(11,4,8)$, respectively. This construction gives the stated degrees
and excludes $4$-cycles.
Base D is obtained from the symplectic generalized quadrangle over
$\mathbb F_5$ by splitting each degree-six point into two degree-three
variables. It has 312 variables, 156 checks, degree pair $(3,6)$,
and girth eight.

For A--D, labels specify the base, lift size, and cyclic ($c$) or
random ($r$) lift.
R1200 uses random socket pairing followed by
degree-preserving switches to remove $4$-cycles.
P1200 and P1600 use degree-constrained progressive edge growth.
The numbers in these R/P labels denote the number of variables.

Q410 and Q412 are $(3,4)$ quasi-cyclic graphs with girths 10 and 12
and circulant sizes 10,112 and 9,088, respectively.
They use the first four columns of Zhang and Wang's exponent
matrix~\cite{zhang2010}, with modified circulant shifts.

Construction details, exponent matrices, and fixed seeds are provided
in the repository linked above. Every realized graph was checked for
simplicity, connectivity, degrees, and girth.

Our KB implementation follows Karimi and
Banihashemi~\cite{karimi2013}, rather than their
software, using exponent-vector messages,
node sums, reverse subtraction, nonzero-message
processing, root retransmission suppression,
and earlier-root deactivation.
Table~\ref{tab:comparison-exp} compares the three methods on $N_g$.
Separate longer-target comparisons in the repository show complementary
performance: KB can be faster on some longer cycle-rich targets.
These implementation-dependent results on synthetic graphs
do not establish superiority across LDPC families.

\subsection{Hybrid accuracy and runtime}
Table~\ref{tab:comparison-exp} compares all three algorithms on the
same target $N_g$, using $\varepsilon=0.08$, $r=101$, and a fixed
interleaving quantum $Q=256$. Here $n$ is the number of variable
nodes and $N_g$ is the exact girth-cycle count.
KB and Exact times are medians of seven runs, while Hybrid times
and outputs average 100 runs. Hybrid time includes both interleaved
procedures and scheduling overhead. MARE is the mean of
$100|\widetilde N_g-N_g|/N_g$ over all 100 outputs,
including exact returns.

\begin{table*}[t]
\centering
\caption{Same-target comparison for $N_g$; running times in milliseconds.
KB and Exact are medians of seven runs; Hybrid and MARE average 100 runs.
MARE is mean absolute relative error in percent.}
\label{tab:comparison-exp}
\footnotesize
\setlength{\tabcolsep}{5pt}
\begin{tabular}{lccrrrrrr}
\hline
Graph & $(d_v,d_c)$ & $g$ & $n$ & $N_g$
& KB & Exact & Hybrid & MARE (\%)\\
\hline
A16c & $(3,6)$ & 6 & $672$ & $2,560$ & 1.507 & 0.497 & 0.799 & 7.7\\
A128r & $(3,6)$ & 6 & $5,376$ & $176$ & 58.701 & 6.215 & 104.750 & 0.0\\
A128c & $(3,6)$ & 6 & $5,376$ & $20,480$ & 14.842 & 5.699 & 1.382 & 8.8\\
B128c & $(3,4)$ & 6 & $2,560$ & $4,608$ & 3.068 & 1.095 & 0.454 & 7.6\\
C128c & $(4,8)$ & 6 & $11,264$ & $204,928$ & 112.565 & 44.423 & 3.196 & 7.6\\
D16c & $(3,6)$ & 8 & $4,992$ & $57,680$ & 75.800 & 18.504 & 4.548 & 6.9\\
R1200 & $(4,8)$ & 6 & $1,200$ & $1,564$ & 47.418 & 4.171 & 24.667 & 6.5\\
P1200 & $(3,6)$ & 8 & $1,200$ & $65$ & 44.640 & 4.211 & 44.721 & 0.0\\
P1600 & $(3,4)$ & 10 & $1,600$ & $11$ & 70.511 & 5.195 & 80.908 & 0.0\\
Q410 & $(3,4)$ & 10 & $40,448$ & $40,448$ & 1125.038 & 128.667 & 90.153 & 7.5\\
Q412 & $(3,4)$ & 12 & $36,352$ & $308,992$ & 2971.170 & 292.054 & 62.125 & 7.6\\
\hline
\end{tabular}
\end{table*}

Zero MARE indicates that exact counting finished first in every run,
as on A128r, P1200, and P1600. Cycle-rich instances can benefit
from sampling, whereas sparse ones may incur interleaving overhead
before exact termination. On R1200, sampling wins all 100 runs,
but the hybrid is slower than standalone exact counting because
state-machine steps have unequal practical costs. At girths 10
and 12, Q410 and Q412 use inverse sampling in every run and
obtain lower hybrid latency than either exact method. These
implementation-dependent results do not establish superiority
across LDPC families.

The observed MARE values around $8\%$ are consistent with a
normal approximation to uncapped inverse sampling. Since $S_r$
has mean $r/p$ and variance $r(1-p)/p^2$, for large $r$ the
relative estimation error is approximately normal with variance
$(1-p)/r$. Hence the expected absolute relative error is
approximately $\sqrt{2(1-p)/(\pi r)}$, or about $7.94\%$ for
$r=101$ and small $p$. This is an asymptotic approximation for
uncapped inverse sampling, not an exact hybrid guarantee;
exact termination can reduce the overall MARE. The guarantee
in Theorem~\ref{thm:hybrid} concerns expected error, so a
100-run empirical MARE may exceed $\varepsilon$.

\section*{AI Statement}
OpenAI ChatGPT (GPT-5.6) assisted with literature discovery, language editing,
and checking mathematical exposition, proofs, and experimental scripts.  The author checked the sources, claims, computations, and final
manuscript and assumes full responsibility for its content.

\begin{thebibliography}{99}

\bibitem{gallager1962}
R. G. Gallager, ``Low-density parity-check codes,''
\emph{IRE Trans. Inf. Theory}, vol. 8, no. 1,
pp. 21--28, 1962.

\bibitem{tanner1981}
R. M. Tanner, ``A recursive approach to low complexity codes,''
\emph{IEEE Trans. Inf. Theory}, vol. 27, no. 5,
pp. 533--547, 1981.

\bibitem{halford2006}
T. R. Halford and K. M. Chugg,
``An algorithm for counting short cycles in bipartite graphs,''
\emph{IEEE Trans. Inf. Theory}, vol. 52, no. 1,
pp. 287--292, 2006.

\bibitem{karimi2013}
M. Karimi and A. H. Banihashemi,
``Message-passing algorithms for counting short cycles in a graph,''
\emph{IEEE Trans. Commun.}, vol. 61, no. 2,
pp. 485--495, 2013.

\bibitem{blake2018}
I. F. Blake and S. Lin,
``On short cycle enumeration in biregular bipartite graphs,''
\emph{IEEE Trans. Inf. Theory}, vol. 64, no. 10,
pp. 6526--6535, 2018.

\bibitem{dehghan2019}
A. Dehghan and A. H. Banihashemi,
``On computing the multiplicity of cycles in bipartite graphs
using the degree distribution and the spectrum of the graph,''
\emph{IEEE Trans. Inf. Theory}, vol. 65, no. 6,
pp. 3778--3789, 2019.

\bibitem{li2015}
J. Li, S. Lin, and K. A. S. Abdel-Ghaffar,
``Improved message-passing algorithm for counting short cycles
in bipartite graphs,''
in \emph{Proc. IEEE Int. Symp. Inf. Theory (ISIT)},
2015, pp. 416--420.

\bibitem{dehghan2020}
A. Dehghan and A. H. Banihashemi,
``Counting short cycles in bipartite graphs:
A fast technique/algorithm and a hardness result,''
\emph{IEEE Trans. Commun.}, vol. 68, no. 3,
pp. 1378--1390, 2020.

\bibitem{chen2025fmpa}
H. Chen and X. He,
``A fast messages-passing algorithm for counting short cycles
in bipartite graphs,''
in \emph{Proc. IEEE/CIC Int. Conf. Commun. China Workshops
(ICCC Workshops)}, 2025, pp. 1--6,
doi: 10.1109/ICCCWorkshops67136.2025.11148122.

\bibitem{alon1997}
N. Alon, R. Yuster, and U. Zwick,
``Finding and counting given length cycles,''
\emph{Algorithmica}, vol. 17, no. 3,
pp. 209--223, 1997.

\bibitem{yang2017}
K. Yang, B. Zhang, Y. Zhan, and D. Guo,
``Design and analysis of efficient algorithm for counting
and enumerating cycles in LDPC codes,''
in \emph{Information Technology and Intelligent Transportation Systems},
Advances in Intelligent Systems and Computing, vol. 454.
Springer, 2017, pp. 131--137,
doi: 10.1007/978-3-319-38789-5\_23.

\bibitem{niu2025}
J. Niu, J. Zola, and A. E. Sar{\i}y\"uce,
``Fast counting and utilizing induced 6-cycles in bipartite networks,''
\emph{IEEE Trans. Knowl. Data Eng.}, vol. 37, no. 6,
pp. 3386--3398, 2025,
doi: 10.1109/TKDE.2025.3546516.

\bibitem{dehghan2020directed}
A. Dehghan and A. H. Banihashemi,
``On computing the number of short cycles in bipartite graphs
using the spectrum of the directed edge matrix,''
\emph{IEEE Trans. Inf. Theory}, vol. 66, no. 10,
pp. 6037--6047, 2020.

\bibitem{fossorier2004}
M. P. C. Fossorier,
``Quasicyclic low-density parity-check codes
from circulant permutation matrices,''
\emph{IEEE Trans. Inf. Theory}, vol. 50, no. 8,
pp. 1788--1793, 2004.

\bibitem{karimi2012qc}
M. Karimi and A. H. Banihashemi,
``Counting short cycles of quasi cyclic protograph LDPC codes,''
\emph{IEEE Commun. Lett.}, vol. 16, no. 3,
pp. 400--403, 2012.

\bibitem{gomez2023}
A. G\'omez-Fonseca, R. Smarandache, and D. G. M. Mitchell,
``An efficient strategy to count cycles in the Tanner graph
of quasi-cyclic LDPC codes,''
\emph{IEEE J. Sel. Areas Inf. Theory}, vol. 4,
pp. 499--513, 2023.

\bibitem{bass1992}
H. Bass,
``The Ihara--Selberg zeta function of a tree lattice,''
\emph{Int. J. Math.}, vol. 3, no. 6,
pp. 717--797, 1992.

\bibitem{mendo2009}
L. Mendo,
``Estimation of a probability with guaranteed normalized
mean absolute error,''
\emph{IEEE Commun. Lett.}, vol. 13, no. 11,
pp. 817--819, 2009.

\bibitem{zhang2010}
G. Zhang and X. Wang,
``Girth-12 quasi-cyclic LDPC codes with consecutive lengths,''
arXiv:1001.3916, 2010.

\end{thebibliography}


\end{document}